% !TeX root = ../main.tex
\chapter*{Preface}
\phantomsection
\addcontentsline{toc}{chapter}{Preface}
\markboth{Preface}{Preface}
\label{ch:preface}

Underground engineering is taught as a sequence of methods and practised as a
sequence of decisions, and the two do not map onto each other neatly. A
graduate can read a $Q$-chart, close a ground--support interaction calculation
and run a finite element model, and still be unable to say which of the three should govern a
particular support decision, or what would have to be true for the answer to
be wrong. This book is organised around that gap.

Its argument is that underground design is a chain, and that the chain is only
as good as the transfer between its links. Site investigation produces
observations; classification converts observations into indices; parameter
derivation converts those indices, with laboratory and in-situ measurement,
into design values; analysis converts design values into predictions; construction tests the predictions; monitoring records the test.
Each conversion introduces assumptions, and each assumption is a place where a
defensible design can quietly become an indefensible one. The book follows
that chain in order, and at each step asks what the method requires, what it
returns, and what it cannot establish.

Two things have changed the chain in the last decade, and both are treated
here as design criteria rather than as commentary. Carbon has become a
quantity that design decisions determine --- through alignment, diameter,
lining thickness and support timing, not principally through material
substitution. Digital delivery has made it feasible to carry ground,
structure, process and performance information in a single model that survives
from feasibility to operation, which is what makes traceability from a
decision back to its evidence practical rather than aspirational. Resilience
is not new, but the performance question it asks has moved from survival to
service continuity, and it is treated here on the same footing. These three
open the book, in Part~\ref{part:context},
because they shape decisions taken before any ground has been characterised.

The remaining parts follow the project. Part~\ref{part:characterisation}
establishes what the ground is: geology, site investigation,
three-dimensional ground modelling, groundwater, in-situ stress, structure and
rock mass classification. Part~\ref{part:parameters} converts that into
parameters and behaviour types. Part~\ref{part:design} designs, running the
three method families in order of increasing resolution and cost --- empirical
for preliminary work, analytical for checking mechanisms, numerical for
detailed verification --- and then treats the support and lining systems that
result, the seismic case, and the reliability and integration framework that
decides between competing designs.
Part~\ref{part:construction} builds. Part~\ref{part:performance} monitors,
operates and maintains, and closes with the risk allocation and contractual
arrangements through which all of it is procured.

The book is written for practising geotechnical and tunnel engineers and for
students approaching the subject with a grounding in soil and rock mechanics.
It assumes familiarity with effective stress, Mohr circles, elementary
elasticity and elementary probability, and introduces what it needs beyond
that where it is first required. Worked calculation sequences are collected
in Appendix~\ref{app:calculations} so that software output can be bounded or checked by hand
wherever an independent answer exists, which remains the most reliable defence against a confident wrong
answer.

Two conventions are stated here because they cause more confusion than any
other. Compression is positive throughout, following rock mechanics practice
rather than structural convention. Principal stresses are ordered
$\sone \geq \stwo \geq \sthree$. Departures from either are flagged
where they occur.

The material grew out of teaching underground design and out of research on
the sustainability and resilience of underground infrastructure. I am grateful to the students whose questions exposed the
weak transfers in the design chain, and to the industry colleagues whose
project experience stands behind the book.

\vspace{\baselineskip}
\noindent\textit{Melbourne, 2026}
